\documentclass[10pt]{article}
\usepackage[paperwidth=36cm, paperheight=35cm, margin=1.5cm]{geometry}
\usepackage[utf8]{inputenc}
\usepackage[T1]{fontenc}
\usepackage[brazil]{babel}

% Uso de fonte sem serifa limpa (se disponível no sistema básico, Latin Modern Sans)
\usepackage{lmodern}
\renewcommand{\familydefault}{\sfdefault}

\usepackage{tikz}
\usetikzlibrary{shapes.geometric, shapes.misc, arrows.meta, positioning, calc, fit, backgrounds, shadows}
\usepackage{xcolor}
\pagestyle{empty}

% Cores Fortes Baseadas no Estilo AHA Original
\definecolor{ahared}{RGB}{205,32,44}      % Vermelho AHA
\definecolor{ahayellow}{RGB}{240,171,0}   % Amarelo Chocável AHA
\definecolor{ahablue}{RGB}{0,111,180}     % Azul RCP AHA
\definecolor{ahagreen}{RGB}{0,133,66}     % Verde Resultados AHA
\definecolor{ahagrey}{RGB}{80,80,80}      % Cinza Escuro para texto e linhas
\definecolor{lightgrey}{RGB}{240,240,240} % Cinza muito claro para caixas laterais

% Estilos TikZ Profissionais
\tikzset{
    basebox/.style={
        rectangle, 
        draw=#1, 
        thick, 
        fill=#1, 
        text=white, 
        rounded corners=4pt, 
        align=left, 
        inner sep=8pt, 
        text width=7.5cm,
        font=\sffamily\normalsize,
        drop shadow={shadow xshift=1.5pt, shadow yshift=-1.5pt}
    },
    redbox/.style={basebox=ahared},
    yellowbox/.style={basebox=ahayellow, text=black},
    bluebox/.style={basebox=ahablue},
    greenbox/.style={basebox=ahagreen},
    decision/.style={
        rectangle, 
        draw=ahagrey, 
        thick,
        fill=white, 
        text=ahagrey,
        rounded corners=2pt, 
        align=center, 
        inner sep=6pt, 
        text width=3.5cm, 
        font=\sffamily\bfseries,
        drop shadow={shadow xshift=1.5pt, shadow yshift=-1.5pt}
    },
    numcirc/.style={
        circle, 
        fill=white, 
        draw=#1, 
        thick,
        text=#1, 
        font=\sffamily\bfseries\large, 
        inner sep=2pt, 
        minimum size=22pt
    },
    line/.style={
        draw=ahagrey, 
        ultra thick, 
        -{Latex[length=3mm, width=2.5mm]}
    },
    sidebox/.style={
        rectangle, 
        draw=ahagrey, 
        thick,
        fill=white, 
        text=black,
        rounded corners=4pt, 
        align=left, 
        inner sep=8pt, 
        text width=8.5cm, 
        font=\sffamily\small,
        drop shadow={shadow xshift=1.5pt, shadow yshift=-1.5pt}
    },
    sideboxTitle/.style={
        fill=ahagrey, 
        text=white, 
        font=\sffamily\bfseries\normalsize, 
        inner xsep=8pt, 
        inner ysep=4pt, 
        align=left
    }
}

% Macros para melhorar listas de itens dentro dos nós sem quebrar espaçamento Tikz
\newcommand{\titem}[1]{\hangindent=3ex\mbox{$\bullet$ }#1\par}

\begin{document}

\begin{center}
\begin{tikzpicture}[auto]

% TÍTULO DO DIAGRAMA
\node[font=\sffamily\huge\bfseries, text=ahared, align=center] (title) at (0, 0.5) {Algoritmo de Parada Cardiorrespiratória Pediátrica (PALS 2025)};

% ==============================
% TRONCO PRINCIPAL
% ==============================
\node[redbox] (box1) at (0, -1.5) {%
    \textbf{Iniciar RCP}\\[1.5ex]
    \titem{Iniciar ventilação com bolsa-máscara e fornecer oxigênio}
    \titem{Instalar monitor/desfibrilador}
};
\node[numcirc=ahared] at ([xshift=-14pt, yshift=14pt]box1.north west) {1};

\node[decision] (shock1) at (0, -4.0) {Ritmo \\ chocável?};
\draw[line] (box1) -- (shock1);

% ==============================
% LADO ESQUERDO: FV/TV sem pulso (x = -5.0)
% ==============================
\node[redbox, align=center] (box2) at (-5.5, -6.0) {\textbf{FV/TV sem pulso}};
\node[numcirc=ahared] at ([xshift=-14pt, yshift=14pt]box2.north west) {2};
\draw[line] (shock1.west) -| node[above, near start, font=\bfseries] {Sim} (box2.north);

\node[yellowbox, align=center] (box3) at (-5.5, -7.8) {\textbf{Choque}\\ \normalsize (o mais rápido possível)};
\node[numcirc=ahayellow] at ([xshift=-14pt, yshift=14pt]box3.north west) {3};
\draw[line] (box2) -- (box3);

\node[bluebox] (box4) at (-5.5, -9.8) {
    \textbf{RCP por 2 min}\\[1.5ex]
    \titem{Acesso IV/IO}
};
\node[numcirc=ahablue] at ([xshift=-14pt, yshift=14pt]box4.north west) {4};
\draw[line] (box3) -- (box4);

\node[decision] (shock2) at (-5.5, -12.3) {Ritmo \\ chocável?};
\draw[line] (box4) -- (shock2);

\node[yellowbox, align=center] (box5) at (-5.5, -14.3) {\textbf{Choque}};
\node[numcirc=ahayellow] at ([xshift=-14pt, yshift=14pt]box5.north west) {5};
\draw[line] (shock2.south) -- node[right, font=\bfseries] {Sim} (box5.north);

\node[bluebox] (box6) at (-5.5, -17.0) {
    \textbf{RCP por 2 min}\\[1.5ex]
    \titem{Epinefrina a cada 3 a 5 min}
    \titem{Considerar via aérea avançada e capnografia}
};
\node[numcirc=ahablue] at ([xshift=-14pt, yshift=14pt]box6.north west) {6};
\draw[line] (box5) -- (box6);

\node[decision] (shock3) at (-5.5, -20.0) {Ritmo \\ chocável?};
\draw[line] (box6) -- (shock3);

\node[yellowbox, align=center] (box7) at (-5.5, -22.0) {\textbf{Choque}};
\node[numcirc=ahayellow] at ([xshift=-14pt, yshift=14pt]box7.north west) {7};
\draw[line] (shock3.south) -- node[right, font=\bfseries] {Sim} (box7.north);

\node[bluebox] (box8) at (-5.5, -24.8) {
    \textbf{RCP por 2 min}\\[1.5ex]
    \titem{Amiodarona ou lidocaína}
    \titem{Tratar causas reversíveis}
};
\node[numcirc=ahablue] at ([xshift=-14pt, yshift=14pt]box8.north west) {8};
\draw[line] (box7) -- (box8);

% Retorno (Loop do choque)
\coordinate (ret5) at (-5.5, -26.8);
\draw[line] (box8.south) -- (ret5) -- ++(-4.5cm,0) |- node[pos=0.25, left, font=\bfseries] {Sim} (box5.west);
\node[font=\sffamily\bfseries\normalsize, text=ahagrey] at (-5.5, -26.5) {Ritmo chocável?};

% ==============================
% LADO DIREITO: Assistolia/AESP (x = 5.0)
% ==============================
\node[redbox, align=center] (box9) at (5.5, -6.0) {\textbf{Assistolia/AESP}};
\node[numcirc=ahared] at ([xshift=-14pt, yshift=14pt]box9.north west) {9};
\draw[line] (shock1.east) -| node[above, near start, font=\bfseries] {Não} (box9.north);

\node[bluebox, align=center] (box10_pre) at (5.5, -7.8) {\textbf{Epinefrina}\\ \normalsize (o mais rápido possível)};
\draw[line] (box9) -- (box10_pre);

\node[bluebox] (box10) at (5.5, -10.3) {
    \textbf{RCP por 2 min}\\[1.5ex]
    \titem{Acesso IV/IO}
    \titem{Epinefrina a cada 3 a 5 min}
    \titem{Considerar via aérea avançada e capnografia}
};
\node[numcirc=ahablue] at ([xshift=-14pt, yshift=14pt]box10.north west) {10};
\draw[line] (box10_pre) -- (box10);

\node[decision] (shock4) at (5.5, -13.5) {Ritmo \\ chocável?};
\draw[line] (box10) -- (shock4);

\node[bluebox] (box11) at (5.5, -16.0) {
    \textbf{RCP por 2 min}\\[1.5ex]
    \titem{Tratar causas reversíveis}
};
\node[numcirc=ahablue] at ([xshift=-14pt, yshift=14pt]box11.north west) {11};
\draw[line] (shock4.south) -- node[right, font=\bfseries] {Não} (box11.north);

\node[decision] (shock5) at (5.5, -19.0) {Ritmo \\ chocável?};
\draw[line] (box11) -- (shock5);

% Cruzamentos da Direita para a Esquerda
\coordinate (retYes10) at (4.0, -13.5);
\coordinate (retYes11) at (4.0, -19.0);

\draw[line] (shock4.west) -- node[above, font=\bfseries] {Sim} (retYes10) |- (box5.east);
\draw[line] (shock5.west) -- node[above, font=\bfseries] {Sim} (retYes11) |- (box7.east);

% ==============================
% CAIXA VERDE (Final)
% ==============================
\node[greenbox, text width=9.5cm] (box12) at (0, -28.5) {
    \titem{Se não houver sinais de retorno da circulação espontânea (RCE), vá para 10 ou 11}
    \vspace{1ex}
    \titem{Se houver RCE, vá para Cuidados Pós-Parada Cardiorrespiratória}
};
\node[numcirc=ahagreen] at ([xshift=-14pt, yshift=14pt]box12.north west) {12};

% Conectar "Não" do Lado Esquerdo para a caixa Verde
\draw[line] (shock2.east) -- node[above, font=\bfseries] {Não} (-1.0, -12.3) -- (-1.0, -28.5) -- (box12.west);
\draw[line] (shock3.east) -- node[above, font=\bfseries] {Não} (-1.0, -20.0);

% Retorno para Green Box (Não do lado esquerdo)
\draw[line] (-5.5, -26.8) -- node[above, font=\bfseries] {Não} (-1.0, -26.8);


% ==============================
% PAINEL LATERAL (Textos Auxiliares)  (x = 12.0)
% ==============================
\node[sidebox] (side1) at (12.0, -1.0) [anchor=north] {
    \textbf{\textcolor{ahared}{Qualidade da RCP}}\\[1.5ex]
    \titem{Comprimir com força ($\ge$ 1/3 do diâmetro AP do tórax) e rápido (100 a 120/min) e permitir o retorno total do tórax.}
    \titem{Minimizar as interrupções nas compressões.}
    \titem{Alternar o responsável pelas compressões a cada 2 min, ou antes se houver cansaço.}
    \titem{Se não houver via aérea avançada, relação compressão-ventilação de 15:2.}
    \titem{Se houver via aérea avançada, fornecer compressões contínuas e 1 ventilação a cada 2 a 3 segundos.}
};

\node[sidebox] (side2) at (12.0, -7.0) [anchor=north] {
    \textbf{\textcolor{ahared}{Energia de Choque}}\\[1.5ex]
    \titem{Primeiro choque 2 J/kg}
    \titem{Segundo choque 4 J/kg}
    \titem{Choques subsequentes $\ge$ 4 J/kg, máx 10 J/kg ou dose de adulto}
};

\node[sidebox] (side3) at (12.0, -10.5) [anchor=north] {
    \textbf{\textcolor{ahared}{Terapia Medicamentosa}}\\[1.5ex]
    \titem{\textbf{Epinefrina:} 0,01 mg/kg IV/IO (0,1 mL/kg da concentração 0,1 mg/mL). Máx 1 mg.}
    \vspace{0.5ex}
    \titem{\textbf{Amiodarona:} 5 mg/kg IV/IO em bolus (máx 300 mg). Repetir até 3 doses. OU}
    \vspace{0.5ex}
    \titem{\textbf{Lidocaína:} 1 mg/kg IV/IO.}
};

\node[sidebox] (side4) at (12.0, -16.0) [anchor=north] {
    \textbf{\textcolor{ahared}{Via Aérea Avançada}}\\[1.5ex]
    \titem{Intubação endotraqueal ou via aérea supraglótica.}
    \titem{Capnografia ou capnometria para confirmar e monitorar.}
};

\node[sidebox] (side5) at (12.0, -19.0) [anchor=north] {
    \textbf{\textcolor{ahared}{Causas Reversíveis}}\\[1.5ex]
    \titem{Hipovolemia \hfill Hipotermia}
    \titem{Hipóxia \hfill Tensão no tórax (pneumotórax)}
    \titem{Hidrogênio (acidose) \hfill Tamponamento cardíaco}
    \titem{Hipoglicemia \hfill Toxinas}
    \titem{Hipo-/hipercalemia \hfill Trombose (pulm./coro.)}
};

\end{tikzpicture}
\end{center}
\end{document}
